\subsection*{4.4 The Split Operation}
In this section we will show how to apply a restriction and then use an operation Split (motivated by [20]) in order to eliminate high quadratic degree terms. Our main focus will be to handle the case where the variable to be set is an extension variable with extension axiom $z-Q=0$ where $z$ is Nonsingular, since in the other case we can potentially just set $z=0$ to eliminate terms. We start by showing how to apply a small Boolean restriction $\rho$ such that $Q|_{\rho}$ is a simple linear function of just one variable.
\paragraph{Lemma22.} Let $z$ be an extension variable with extension axiom $z-Q(w_1,\ldots,w_k)$, for $k\leq\kappa$. Assume that $z$ is Nonsingular (i.e. $\supp(Q)\subseteq[1,\ldots,p-1]$) and $|\supp(Q)|\geq2$. Then for every $l\in[0,\ldots,\ell(\supp(Q))-1]$, there exists a variable $w^*$ in $\vars(Q)$, and a restriction $\delta$ to $\vars(Q)-w^*$ such that:
\begin{enumerate}[label=(\arabic*)]
\item $Q|_{\delta}=(b-a)w^*+a$, where $b,a\in\supp(Q)$. Thus $Q|_{\delta}$ is a linear function of $w^*$ and $\supp(Q|_{\delta})=\{a,b\}$;
\item $a^l\neq b^l\pmod p$.
\end{enumerate}
\begin{proof}
We will create a decision tree that will query $\vars(Q)$ one-by-one. Associated with the root $r$ is the set of values $S_r=\{a^l\mid a\in\supp(Q)\}$. That is, we label the root with the set of all possible values that $z^l$ can take on. Since $l<\ell(\supp(z))$, it follows that $|S_r|\geq2$ (since otherwise we would have $l=\ell(\supp(z))$). At the root we query the first variable $w_1$, with left edge labelled by $w_1=0$ and right edge labelled by $w_1=1$. Now we label the left vertex with the set $\{a^l\mid a\in\supp(Q|_{w_1=0})\}$, of all values that $z^l$ can take on under the restriction $w_1=0$. Similarly we label the right vertex with the set $\{a^l\mid a\in\supp(Q|_{w_1=1})\}$. We continue recursively, querying the next variable at each vertex $v$ of the decision tree, as long as the set of allowable values for $z^l$ under the partial restriction $\rho_v$ associated with $v$ is greater than one. Now consider the longest path, $\xi$ in $T$. The partial restriction $\rho$ associated with $\xi$ sets the first $k'$ variables, where $k'\geq1$ since initially $z^l$ takes on at least two values. Also since $\xi$ is a complete path, the associated set $\{a^l\mid a\in\supp(Q|_{\rho})\}$ contains exactly one element, call it $q$.
Now consider the twin path $\xi'$ with associated restriction $\rho'$, where $\rho'$ is obtained from $\rho$ by toggling the value of the last variable, $w_{k'}$, queried. Again since $\xi'$ is a complete path, the associated set $\{a^l\mid a\in\supp(Q|_{\rho'})\}$ contains exactly one element, call it $q'$. Note that $q,q'$ must be distinct.
Let $\delta$ be the following assignment to $\vars(Q)-w_{k'}$: for $1\leq j<k'$, we set $\delta(w_j)=\rho(w_j)=\rho'(w_j)$, and for $k'<j\leq k$, we set $\delta(w_j)=0$. Setting $w^*=w_{k'}$, $Q|_{\delta}$ is a linear equation of the form $(b-a)w^*+a$, where $b,a\in\supp(Q)$. Finally, by construction, $a^l\neq b^l$ (since otherwise the two paths corresponding to $\rho,\rho'$ would be the same).
\end{proof}
In the remainder of this subsection, we will be interested in the case where we want to eliminate some Nonsingular extension variable $z$ from the refutation, and we have already applied the above Lemma so that the extension axiom for $z$ is of the form $z-((b-a)w+a)=0$, where $w$ is some variable in $\mathbf x\cup\mathbf y$. Thus, $\supp(z)=\{a,b\}$. The next two Lemmas generalizes a similar argument due to Sokolov, and show how to remove Quadratic degree pairs of the form $(t_1z^i,t_2z^j)$ for a carefully chosen pair $i,j$ from the refutation via the Split operation.
\paragraph{Lemma23.} Let $z$ be an extension variable such that $\supp(z)=\{a,b\}$, where $a\neq b$ and $a,b\in\mathbb F_p^*$ and let $P$ be any polynomial. Then, for any two distinct numbers $\ell_0,\ell_1$ where $\ell_0<\ell_1$ and $a^{\ell_1-\ell_0}\neq b^{\ell_1-\ell_0}$, there exists a unique polynomial $R=R_0z^{\ell_0}+R_1z^{\ell_1}$ such that $R=P\bmod(z-a)(z-b)$. That is, $R(a)=P(a)$ and $R(b)=P(b)$, where $P(a)$ denotes the polynomial $P$ under the substitution $z=a$.
\begin{proof}
Let $z-Q=0$ be the extension axiom for $z$, where $\supp(z)=\{a,b\}$. Then by Lemma13 the polynomial $(z-a)(z-b)=0$ is implied by (and derivable from) the extension axiom for $z$ plus the Boolean axioms. We can assume without loss of generality that $P$ has the form $P_0+zP_1+\ldots+z^{p-2}P_{p-2}$.
Now we want to argue that there exists a polynomial $R=z^{\ell_0}R_0+z^{\ell_1}R_1$, where $R_0,R_1$ are polynomials over $\vars(P)-z$, and such that $R(a)=P(a)$, and $R(b)=P(b)$. We can find $R_0$ and $R_1$ by solving the following system of equations, where we view $R_0,R_1$ as the underlying variables, and treating $P(a),P(b)$ as constants:
\[a^{\ell_0}R_0+a^{\ell_1}R_1=P(a),\qquad b^{\ell_0}R_0+b^{\ell_1}R_1=P(b).\]
This has a (unique) solution since the determinant of the associated matrix is
\[\begin{vmatrix}a^{\ell_0}&a^{\ell_1}\\b^{\ell_0}&b^{\ell_1}\end{vmatrix}=a^{\ell_0}b^{\ell_0}(b^{\ell_1-\ell_0}-a^{\ell_1-\ell_0}).\]
By our assumption, this matrix is non-singular over $\mathbb F_p$ and therefore the above system of equations has a unique solution over $\mathbb F_p$, given by:
\[\binom{R_0}{R_1}=\begin{pmatrix}a^{\ell_0}&a^{\ell_1}\\b^{\ell_0}&b^{\ell_1}\end{pmatrix}^{-1}\binom{P(a)}{P(b)}.\]
Abbreviating $a^{\ell_0},a^{\ell_1},b^{\ell_0},b^{\ell_1}$ by $a_0,a_1,b_0,b_1$ respectively, we have by definition of the inverse:
\[\binom{R_0}{R_1}=\begin{pmatrix}a_0&a_1\\b_0&b_1\end{pmatrix}^{-1}\binom{P(a)}{P(b)}=\frac{1}{a_0b_1-a_1b_0}\begin{pmatrix}b_1&-a_1\\-b_0&a_0\end{pmatrix}\binom{P(a)}{P(b)}.\]
Solving for $R_0$ we have:
\begin{align*}
R_0&=\frac{b_1}{a_0b_1-a_1b_0}P(a)-\frac{a_1}{a_0b_1-a_1b_0}P(b)\\
&=\frac{b_1}{a_0b_1-a_1b_0}\left(a_0P_{\ell_0}+a_1P_{\ell_1}+\sum_{i\neq\ell_0,\ell_1}a^iP_i\right)\\
&\quad-\frac{a_1}{a_0b_1-a_1b_0}\left(b_0P_{\ell_0}+b_1P_{\ell_1}+\sum_{i\neq\ell_0,\ell_1}b^iP_i\right)\\
&=\frac{a_0b_1}{a_0b_1-a_1b_0}P_{\ell_0}+\frac{a_1b_1}{a_0b_1-a_1b_0}P_{\ell_1}\\
&\quad+\frac{b_1}{a_0b_1-a_1b_0}\left(\sum_{i\neq\ell_0,\ell_1}a^iP_i\right)-\frac{a_1b_0}{a_0b_1-a_1b_0}P_{\ell_0}\\
&\quad-\frac{a_1b_1}{a_0b_1-a_1b_0}P_{\ell_1}-\frac{b_1}{a_0b_1-a_1b_0}\left(\sum_{i\neq\ell_0,\ell_1}b^iP_i\right)\\
&=P_{\ell_0}+\sum_{i\neq\ell_0,\ell_1}c_{0i}P_i
\end{align*}
for some constants $c_{0i}\in\mathbb F_p$. And similarly solving for $R_1$, it has the following form:
\[R_1=P_{\ell_1}+\sum_{i\neq\ell_0,\ell_1}c_{1i}P_i\]
for some constants $c_{1i}\in\mathbb F_p$.
\end{proof}
\paragraph{Definition24 (Split).} Let $z$ be an extension variable with extension axiom $z-Q=0$ such that $\supp(z)=\{a,b\}\subseteq[1,\ldots,p-1]$. For any polynomial $P$ and for every $\ell_0<\ell_1$ such that $a^{\ell_1-\ell_0}\neq b^{\ell_1-\ell_0}$, let $R=R_0z^{\ell_0}+R_1z^{\ell_1}$ be the unique polynomial given by Lemma23 such that $R=P\bmod(z-a)(z-b)$. Then $\Split_{z,\ell_1,\ell_0}(P)$ is defined to be the pair of polynomials $\{R_0,R_1\}$. For a proof $\Pi$, and an extension variable $z$ such that $\supp(z)=\{a,b\}$, we define $\Split_{z,\ell_0,\ell_1}(\Pi)$ to be the sequence of lines $\Split_{z,\ell_0,\ell_1}(P)$, over all $P\in\Pi$.
\paragraph{Lemma25.} Let $\Pi$ be a refutation of a set of unsatisfiable polynomials $F$. Let $z$ be a variable that occurs in $\Pi$ such that the polynomials in $F$ do not contain $z$ except for the axiom $(z-a)(z-b)=0$ for some $a,b\in\mathbb F_p^*$. Then for any $\ell_0,\ell_1$ such that $\ell_0<\ell_1$ and $a^{\ell_1-\ell_0}\neq b^{\ell_1-\ell_0}$, $\Pi'=\Split_{z,\ell_0,\ell_1}(\Pi)$ forms a valid refutation of $F$ modulo $(z-a)(z-b)$.
\begin{proof}
Fix an extension variable $z$ in $\Pi$ such that it does not occur in any axioms except $(z-a)(z-b)=0$, and let $\ell_0,\ell_1$ be such that $\ell_0<\ell_1$ and $a^{\ell_1-\ell_0}\neq b^{\ell_1-\ell_0}$. We will show by induction on the number of lines in $\Pi$ that $\Split_{z,\ell_0,\ell_1}(\Pi)$ is a valid derivation that meets the conditions of the lemma. For the base case, note that all of the axioms are either free of $z$ or eliminated as a result of reducing by $(z-a)(z-b)$, and hence their Split versions are derivable. Now suppose that the Lemma holds for the first $j-1$ lines of $\Pi$; that is, $\Split_{z,\ell_0,\ell_1}(\Pi_{j-1})$ is a valid derivation, where $\Pi_{j-1}$ denotes the first $j-1$ lines of $\Pi$.
The first case is where $P_j$ is a linear combination of two previously derived lines, so $P_j=\alpha P_{j_1}+\beta P_{j_2}$ for some $j_1$ and $j_2$ less than $j$ and $\alpha,\beta\in\mathbb F_p$. Using the inductive hypothesis, we have:
\begin{align*}
P_j&=\alpha(z^{\ell_0}R_{j_10}+z^{\ell_1}R_{j_11})+\beta(z^{\ell_0}R_{j_20}+z^{\ell_1}R_{j_21})\pmod{(z-a)(z-b)}\\
&=z^{\ell_0}(\alpha R_{j_10}+\beta R_{j_20})+z^{\ell_1}(\alpha R_{j_11}+\beta R_{j_21})\pmod{(z-a)(z-b)}.
\end{align*}
By the uniqueness of the polynomial $R_j=z^{\ell_0}R_{j0}+z^{\ell_1}R_{j1}$ that is equivalent to $P_j\bmod(z-a)(z-b)$ (by Lemma23), this implies that $R_{j0}=\alpha R_{j_10}+\beta R_{j_20}$ and similarly $R_{j1}=\alpha R_{j_11}+\beta R_{j_21}$, and thus $R_{j0}$ can be derived from a linear combination of $R_{j_10}$ and $R_{j_20}$ and similarly for $R_{j1}$.
The second case is when $P_j$ is derived from a previously derived line $P_{j'}$ by multiplying $P_{j'}$ by a variable $w$. That is, $P_j=wP_{j'}$ for some $j'<j$. If $w\neq z$, then we have that $R_{j1}=wR_{j'1}$ (similarly for $R_{j0}$). If $w=z$ then we have:
\[\binom{R_{j'1}}{R_{j'0}}=\begin{pmatrix}a^{\ell_1}&a^{\ell_0}\\b^{\ell_1}&b^{\ell_0}\end{pmatrix}^{-1}\binom{P_{j'}(a)}{P_{j'}(b)}\]
from which we need to derive
\begin{align*}
\binom{R_{j1}}{R_{j0}}&=\begin{pmatrix}a^{\ell_1}&a^{\ell_0}\\b^{\ell_1}&b^{\ell_0}\end{pmatrix}^{-1}\binom{P_j(a)}{P_j(b)}\\
&=\begin{pmatrix}a^{\ell_1}&a^{\ell_0}\\b^{\ell_1}&b^{\ell_0}\end{pmatrix}^{-1}\binom{aP_{j'}(a)}{bP_{j'}(b)}\\
&=\begin{pmatrix}a^{\ell_1}&a^{\ell_0}\\b^{\ell_1}&b^{\ell_0}\end{pmatrix}^{-1}\begin{pmatrix}a&0\\0&b\end{pmatrix}\begin{pmatrix}a^{\ell_1}&a^{\ell_0}\\b^{\ell_1}&b^{\ell_0}\end{pmatrix}\binom{R_{j'1}}{R_{j'0}}.
\end{align*}
Thus, $R_{j1}$ can be derived as a linear combination of $R_{j'1}$ and $R_{j'0}$, and similarly for $R_{j0}$.
\end{proof}
